\section{The random Helly-size study (main Remark~\mref{rem:random})}\label{si:random}

\paragraph{Protocol.} For each dimension $k$, $m=k+3$ rows $\psi_a(u)=c_a+v_a^\top u$ are drawn with independent integer coefficients uniform in $\{-3,\dots,3\}$ (both $c_a$ and the entries of $v_a$), on the action set $U=[-1,1]^k$; the pseudo-random seed is fixed in the script. A system is \emph{infeasible} when its lower value $\max_{u\in U}\min_a\psi_a(u)$ is negative; for every infeasible system the smallest $s$ such that some $s$-subfamily is infeasible is located by enumerating subfamilies in order of increasing size. For $k\le3$ every feasibility test is the exact lower-value programme (vertex enumeration of the lifted polytope in $(u,t)$). For $k=4$ the tests use a float screen: HiGHS proposes a maximiser, the maximiser is rationalised and checked exactly against every row and every facet of $U$ (a rational point that passes is an exact feasibility certificate); whenever the screen does not produce an exact certificate the exact programme decides. Check S3.6 confirms that on $k=3$ the screened and the fully exact protocols return identical histograms. In the rank-one variant all coefficient vectors are integer multiples of one random vector $g\in\{-3,\dots,3\}^3$.

\begin{table}[H]
\centering\footnotesize
\caption{Minimal infeasible subfamily sizes among the infeasible systems drawn (checks S3.1--S3.5).}
\begin{tabular}{lccccccccl}
\toprule
case & $m$ & draws & infeasible & size 1 & size 2 & size 3 & size 4 & size 5 & tests\\
\midrule
$k=1$ & 4 & 300 & 217 & 163 & 54 & 0 & 0 & 0 & exact\\
$k=2$ & 5 & 300 & 214 & 89 & 116 & 9 & 0 & 0 & exact\\
$k=3$ & 6 & 150 & 86 & 12 & 52 & 20 & 2 & 0 & exact\\
$k=3$ (rank-one) & 6 & 100 & 85 & 38 & 47 & 0 & 0 & 0 & exact\\
$k=4$ & 7 & 120 & 65 & 2 & 29 & 29 & 5 & 0 & float screen + exact\\

\bottomrule
\end{tabular}
\end{table}

\paragraph{Reading.} The Helly bound $k+1$ is attained in every dimension where enough systems were drawn ($k\le3$) but with decreasing frequency; for $k=4$ no size-$5$ witness appeared among $65$ infeasible systems, while main Theorem~\mref{thm:helly}(ii) shows that such systems exist. In the rank-one variant the bound of the theorem is $r+1=2$ and the observed maximum is $2$. The frequencies depend on the sampling law (integer coefficients in a small range produce many single-row infeasibilities, i.e. rows negative on all of $U$) and are reported as observations, not as a claim about typical certificates in applications. The search cost of this protocol---$\sum_{j\le s}\binom{m}{j}$ feasibility tests per system---is what main Remark~\mref{rem:search} contrasts with the single linear programme that returns a maximal-margin certificate directly.
